\documentclass[11pt,a4paper]{article}
\usepackage[hmargin=26mm,vmargin=24mm]{geometry}
\usepackage[T1]{fontenc}
\usepackage{amsmath,amssymb,amsthm,microtype}
\usepackage[dvipsnames]{xcolor}
\usepackage[colorlinks=true,linkcolor=MidnightBlue,citecolor=MidnightBlue,urlcolor=MidnightBlue]{hyperref}
\hypersetup{pdftitle={Pointwise representation counts for Eisenstein restricted sums},pdfauthor={Roy van Rijn}}
\newtheorem{theorem}{Theorem}[section]
\newtheorem{lemma}[theorem]{Lemma}
\newtheorem{corollary}[theorem]{Corollary}
\newcommand{\A}{\mathcal A}
\newcommand{\Oq}{\mathcal O}
\newcommand{\Z}{\mathbb Z}
\newcommand{\Q}{\mathbb Q}
\newcommand{\C}{\mathbb C}
\newcommand{\F}{\mathbb F}
\newcommand{\nrd}{\operatorname{nrd}}
\newcommand{\PGL}{\operatorname{PGL}}
\newcommand{\SU}{\operatorname{SU}}
\newcommand{\Sym}{\operatorname{Sym}}
\title{Pointwise representation counts for Eisenstein restricted sums}
\author{Roy van Rijn}
\date{}
\begin{document}
\maketitle
\begin{abstract}
Let \(A\) consist of the positive integers supported on primes
\(1\bmod 3\), including \(1\). We count distinct ordered representations
\(n=a+b\) for \(n\equiv2\bmod6\) and \(n=a+3b\) for \(n\equiv4\bmod6\),
with \(a,b\in A\). In either family prescribe an interval of length
\(H\) centered at \(x\) for one actual summand. Put \(M=\min(x,n-x)\).
For every fixed \(0<\varepsilon<1\), if \(H\le M\) and
\(H^3\ge n^{1+\varepsilon}M\), the count has order \(HJ(n)/\log n\)
for every sufficiently large admissible \(n\). Here \(J(n)\) is the
product of \((p-1)/(p-2)\) over the odd primes \(p\equiv2\bmod3\)
dividing \(n\). In particular, central windows of width \(n^\theta\)
are allowed for \(2/3<\theta<1\), and windows at a prescribed small
summand of size comparable to their width are allowed for
\(1/2<\theta<1\). The lower bound combines angular Hecke estimates
with a sieve that prescribes small split factors in both summands.
\end{abstract}
\section{The pointwise count}
Let \(\A=\{a\in\Z_{>0}:p\mid a\Rightarrow p\equiv1\bmod3\}\), where \(p\)
always denotes a prime. Write \(3\!\cdot\!\A=\{3a:a\in\A\}\).
We count ordered pairs \((s_1,s_2)\) of positive integers summing to
\(n\), with
\[
(s_1,s_2)\in
\begin{cases}
\A\times\A,&n\equiv2\pmod6,\\
\A\times(3\!\cdot\!\A)\text{ or }(3\!\cdot\!\A)\times\A,
 &n\equiv4\pmod6.
\end{cases}
\]
In the second line one orientation is prescribed in advance.
For \(0<x<n\) let \(r(n;x,H)\) count these pairs with
\(|s_1-x|\le H/2\). In particular, if the first summand is \(3a\),
it is \(3a\) that lies in the interval.

\begin{theorem}\label{thm:main}
Fix \(0<\varepsilon<1\) and put \(M=\min(x,n-x)\). Uniformly for
every sufficiently large admissible \(n\) and every real \(x,H\) with
\[
0<x<n,\qquad 1\le H\le M,\qquad H^3\ge n^{1+\varepsilon}M,
\]
one has, in each of the prescribed families,
\begin{equation}\label{eq:main}
r(n;x,H)\asymp_\varepsilon \frac{HJ(n)}{\log n},
\qquad
J(n)=\prod_{\substack{p\mid n,\ p\ge5\\p\equiv2\ (3)}}
             \frac{p-1}{p-2}.
\end{equation}
The constants in the two bounds and the threshold for \(n\) depend
only on \(\varepsilon\).
\end{theorem}

The product has a direct interpretation. At \(p=5\), for example,
the two conditions \(5\nmid s_1\) and \(5\nmid n-s_1\) forbid two
residues when \(5\nmid n\), and one when \(5\mid n\). The ratio of
the allowed proportions is \((4/5)/(3/5)=4/3\). The same calculation
explains every factor of \(J(n)\) and yields an upper bound by an
ordinary sieve on integers.

The lower bound begins with a larger family that can be counted
spectrally. Write \(N(z)=z\bar z\) for the Eisenstein norm on
\(\Z[\omega]\), with \(\omega^2+\omega+1=0\). The equation
\(N(z)+N(w)=n\) is a quaternion norm equation. At an inert prime
\(p\equiv2\bmod3\), \(p\mid N(z)\) forces \(z\in p\Z[\omega]\).
It is therefore possible to remove the large inert primes by a
lattice count within the prescribed window.

A separate issue is that many coordinate pairs can give the same
two norm values. Sifting out small split primes bounds this
multiplicity, with a baseline coordinate count of size
\(H/(\log n)^2\) before the target factors. To gain a factor of \(\log n\)
we allow small split primes to occur to exponent one in one
summand, or in both summands when they divide the target.
For instance, \(7\nmid n\) permits placing \(7\) in either summand.
If \(7\mid n\), its use as a common factor reduces the target to
\(n/7\). The split-prime local factors for these two situations
supply the same logarithmic gain. The prescribed factors can
be read off the resulting values, and each pair then has bounded
weighted coordinate multiplicity.

The angular estimate costs \(\sqrt{M/H}\) in the Hecke remainder.
Multiplying by the square-root error from the norm shell gives
\(\sqrt{nM/H}\). Comparing this with the main-term width \(H\)
explains \(H^3\ge n^{1+\varepsilon}M\). The proof below treats
the finite level and the angular variable together, which lets
us prescribe the mixed orientation at any center.

\paragraph{Proof overview.}
Section~\ref{sec:angular} estimates the angularly weighted shell
with finite congruence tests. Section~\ref{sec:local} handles
primes dividing the target, with a \hyperref[sec:split-example]{worked example at \(7^2\)}.
Section~\ref{sec:sieve} sums over the small split factors to
recover the logarithm. Section~\ref{sec:count} removes the
remaining inert primes, bounds the weight of each distinct
pair, and gives the matching integer-sieve upper bound.

\paragraph{Relation to earlier work.}
Dietmann and Elsholtz \cite{DietmannElsholtz} prove eventual
binary existence for the Gaussian prime restriction \(p\equiv1\bmod4\)
and obtain \(\gg_\delta n^{1-\delta}\) pairs for every \(\delta>0\)
without restricting the summand sizes. Howe, Huffman, Kohlmeyer and
Kruer \cite[Theorem~1.1]{Euler} prove existence for every positive
target \(2\bmod4\), using quaternionic norm counts and a sieve on the
representing vectors. The prime sets used here and in the Gaussian
problem are incomparable, as \(7\) and \(17\) show.
We use the same norm-counting principle for the Eisenstein ring.
The new estimates needed here are angular localization with prescribed
residue conditions, the densities when a sieving prime divides the
target, and a lower bound that counts distinct value pairs at the
integer-sieve scale. Ramar\'e and Viswanadham
\cite[Introduction and Theorem~1.6]{RV} discuss the split-prime
Eisenstein set in their sieve framework and prove a ternary
additive theorem for primitive Gaussian values.

The Frobenian correlation theorem of Loughran and Matthiesen
\cite[Theorem~1.9]{LM} requires pairwise independent linear parts.
Fixing the total gives the one-variable forms \(t,n-t\), with
proportional linear parts. In the mixed case, restricting to the
necessary residue class and dividing one form by \(3\) leaves
the same dependence. Freeing the total gives a two-variable
sum, which averages over the slices that our theorem estimates
individually.

Iwaniec's half-dimensional sieve \cite[Theorem~1]{IwaniecHalf}
uses a half-dimensional local product and a remainder from
counts in residue classes. Sieving both summands directly
excludes two residues per inert prime away from the target
divisors. As the inert primes have density \(1/2\), this gives
sieve dimension one. A half-dimensional sieve on one summand
with the other already restricted would require progression
estimates for that restricted sequence in translated windows,
uniformly in the target. Greaves \cite[main theorem]{Greaves}
obtains a bound of order at least \(n/(\log n)^{5/2}\) for
prime pairs whose complementary value is a sum of two squares.
The cutoff \(p,q\le\sqrt n/2\) in that Gaussian input leaves the
split of the complementary sum of two squares free. Here the
angular Hecke estimate supplies the window-dependent distribution,
and the split-factor sum reaches the scale of distinct values.

\begin{corollary}\label{cor:ranges}
The count in every fixed interval of positive length for \(s_1/n\)
inside \((0,1)\) is of order \(nJ(n)/\log n\). In particular the
unrestricted count has this order. At any fixed interior share,
windows of width \(n^\theta\) have count of order
\(n^\theta J(n)/\log n\) for \(2/3<\theta<1\).
For \(1/2<\theta<1\) and fixed \(0\le a<b\), the same count holds
with \(s_1\in[a n^\theta,b n^\theta]\).
\end{corollary}

The results give two-sided order bounds. A leading constant and
a numerical onset are open here, as are the central and
small-summand endpoint exponents \(2/3\) and \(1/2\).

\section{An angular estimate on a quaternion shell}\label{sec:angular}
Fix \(n,x,H\) as in Theorem~\ref{thm:main}. We first count coordinate
representations with an angular weight and prescribed local conditions.

\subsection{The norm shell}
Let \(R=\Z[\omega]\) and \(K=\Q(\omega)\). The Eisenstein lattice is
norm Euclidean, since every complex number is within squared
distance \(1/3\) of a lattice point. Write \(B=K+Kj\), where
\(j^2=-1\) and \(jz=\bar z j\), and put \(\Oq=R+Rj\).
Its reduced norm is \(\nrd(z+wj)=N(z)+N(w)\). The trace pairing on
\(1,\omega,j,\omega j\) has determinant \(-9\), so the order has
reduced discriminant \(3\). The algebra is definite and the number
of ramified places is even, so it is ramified at \(3\) and infinity.
Thus \(\Oq\) is maximal.

Rounding both complex coordinates gives a remainder of reduced norm
at most \(2/3\). In any nonzero integral right ideal, take a least
positive norm element \(\alpha\) and divide on the left by \(\alpha\).
The remainder in that ideal has norm smaller than \(\nrd(\alpha)\)
unless it vanishes. Every right ideal is therefore principal, as
are the fractional ideals by scaling. The unit group \(U\) consists
of the twelve elements of \(R^\times\cup R^\times j\). For
\(3\nmid m\), let \(\mathcal Q_m=\{\alpha\in\Oq:\nrd(\alpha)=m\}\).
At each \(\ell^e\Vert m\), integral right ideals correspond locally
to index-\(\ell^e\) sublattices of \(\Z_\ell^2\), of which there are
\(1+\ell+\cdots+\ell^e\). Each ideal has twelve generators of norm
\(m\), giving
\begin{equation}\label{eq:shell}
|\mathcal Q_m|=12\sigma(m),\qquad \sigma(m)=\sum_{d\mid m}d .
\end{equation}
We use this all-ideal convention throughout
\cite[Propositions~41.3.1 and~41.3.6]{Voight}.

The number of coordinates of a positive norm \(a\) is \(6f(a)\),
where \(f\) is multiplicative. Its local values are \(f(p^e)=e+1\)
at split primes \(p\equiv1\bmod3\), \(f(p^e)=1\) at \(3\), and
\(f(p^e)=\mathbf1_{2\mid e}\) at inert primes \(p\equiv2\bmod3\).
Indeed a split prime has two conjugate factors in \(R\), and a choice
of their exponents gives each possibility up to a unit. The ramified
prime in \(R\) above \(3\) is generated by \(1-\omega\) and has norm
\(3\). In particular, \(f(a)\le\tau(a)\), where \(\tau\) is the
divisor function. At an inert prime, \(R/pR\) is a field, so norm
divisibility by \(p\) is equivalent to coordinate divisibility by
\(p\). For even \(m\), both component norms have the same parity.
Both are even precisely on the shell scaled by \(2\) from \(m/4\).
The odd-component shell therefore has size
\begin{equation}\label{eq:odd-shell}
 R_o(m)=12\{\sigma(m)-\sigma(m/4)\},\qquad
 \frac{R_o(m)}m=18\prod_{\substack{p^e\Vert m\\p\ge5}}
                            \frac{\sigma(p^e)}{p^e},
\end{equation}
where \(\sigma(m/4)=0\) if \(4\nmid m\).

\subsection{Finite tests and the Hecke bound}
Use the finite group \(G\) obtained by multiplying the fixed factor
\(G_3=(\Oq/3\Oq)^\times/\F_3^\times\) and any factors
\(\PGL_2(\Z/p^a\Z)\) at \(p\ge5\), with \(a=1\) or \(2\).
The Hecke index in this subsection is coprime to the level.
Write \(\Gamma\) for the image of \(U\) in \(G\).
The kernel on \(U\) is \(\{\pm1\}\), as is seen at \(3\).
Each group has a norm squareclass map, using
determinant at split quaternion places. Both squareclasses occur.
At \(3\), for example, the norms of \(1\) and \(1+j\) give the two
classes. They are detected modulo \(p\) at every odd level prime.

For \(k\ge0\), let \(\rho_k=\Sym^{2k}\C^2\) denote the unitary
representation of \(\SU(2)\) on that space. For a real quaternion
of nonzero norm use \(u_\alpha=\alpha/\sqrt{\nrd(\alpha)}\in\SU(2)\),
under the usual identification of unit quaternions with
\(\SU(2)\). The action on order units descends to \(\Gamma\).
Use the space of functions \(F:G\to\Sym^{2k}\C^2\) satisfying
\(F(\gamma g)=\rho_k(\gamma)F(g)\), with squared norm summed once
per left coset. The all-ideal Hecke operator is
\begin{equation}\label{eq:hecke}
 (T_m F)(g)=\frac1{12}\sum_{\alpha\in\mathcal Q_m}
                 \rho_k(u_\alpha)^{-1} F(\alpha_{\rm fin} g),
\end{equation}
where \(\alpha_{\rm fin}\) is projective reduction in \(G\).
Let \(\mathcal K\) be the product of the projective order-unit
groups at the finite places, and \(\mathcal K_0\) the kernel of
reduction to \(G\). At \(3\) we use order units, with reduction
modulo \(3\Oq_3\). A residue unit lifts to an order unit since its
reduced norm is a \(3\)-adic unit. At each \(p\ge5\) in the level,
the corresponding reduction of \(\mathrm{GL}_2(\Z_p)\) is
surjective modulo \(p^a\). Thus \(\mathcal K/\mathcal K_0=G\).

Class number one gives
\(PB^\times(\mathbb A_f)=PB^\times(\Q)\mathcal K\), where
\(\mathbb A_f\) denotes the finite adeles. A rational element
stabilizing \(\mathcal K\) has a representative locally equal to
a scalar times an order unit. A single rational scalar cancels
the valuations of the local scalars, giving an element of \(U\).
Consequently the finite adelic quotient is \(\Gamma\backslash G\).
The projective rational stabilizer at any level coset is trivial.
Indeed, \(\mathcal K_0\) is normal in \(\mathcal K\) and only the
units \(\pm1\) lie in the kernel of reduction. With coefficients,
rational left equivariance gives the stated action of \(\Gamma\)
\cite[Section~2, (4)--(6) and Proposition~2.5]{GV}.
Each orbit has \(|\Gamma|\) elements and hence the same finite
measure. At infinity Schur orthogonality for a unit test vector
gives the factor \(1/(2k+1)\). The resulting automorphic norm
differs from the counting norm by a common scalar.

To see the direction and normalization in \eqref{eq:hecke},
lift \(g\) to \(\tilde g\in\mathcal K\) supported at level primes.
A right-ideal representative \(a\) in the correspondence of norm
\(m\) can be chosen supported at primes of \(m\), so
\(\tilde g a=a\tilde g\). If the ideal is \(\alpha\Oq\), write
\(a=\alpha u\) with \(u\in\mathcal K\) projectively. At level primes
\(u\) reduces to \(\alpha_{\rm fin}^{-1}\). Hence the summand is
\(\rho_k(u_\alpha)F(\alpha_{\rm fin}^{-1}g)\).
Changing \(\alpha\) by a right order unit cancels by equivariance.
Averaging over all twelve generators implements the sum once per
ideal. Replacing \(\alpha\) by \(\bar\alpha\) gives
\eqref{eq:hecke}. Conjugation also proves self-adjointness.

We use the bound
\begin{equation}\label{eq:hecke-norm}
 \|T_m\|\le \tau(m)\sqrt m
\end{equation}
on the whole space for \(k>0\), and on the complement of the joint
norm-squareclass characters for \(k=0\). Here are the relevant
hypotheses for cuspidal transfer. A finite-dimensional smooth
irreducible local factor at a split place \(\ell\) is trivial on
the norm-one group, since its open normal kernel contains the
elementary unipotents by diagonal conjugation. Such a factor
cannot occur at positive weight. In fact \(\Oq[1/\ell]^1\) has
dense real image in \(\SU(2)\). Equation \eqref{eq:shell} gives
infinitely many norm-one elements \(\alpha/\ell^a\) of exact
denominator \(\ell^a\) by taking \(\nrd(\alpha)=\ell^{2a}\) and
omitting \(\ell\Oq\). There are at most \(12(2a+1)\) single-component
points of that norm, so these elements extend outside the Eisenstein
torus normalizer. A proper infinite closed subgroup of \(\SU(2)\)
normalizes its one-dimensional identity component. It contains the
order-three unit \(\omega\), which must then lie in that torus and
determines it as the Eisenstein torus. This contradicts the points
outside its normalizer. For a finite adelic representative \(g\),
impose the congruences \(g^{-1}\gamma g\) belonging to the level
away from \(\ell\). They cut out a finite-index subgroup of
\(\Oq[1/\ell]^1\), still dense by connectedness. At \(\ell\) the
discrepancy is absorbed by invariance under the local norm-one
group. Rational left invariance then forces the coefficient at
\(g\) to be fixed by that dense subgroup, which is impossible
in the nontrivial irreducible representation.

For \(k=0\), suppose the factor at \(\ell\) is finite-dimensional
and let \(h\) be a finite adelic norm-one element.
Strong approximation for \(B^1\) away from \(\{\infty,\ell\}\)
lets us choose a rational norm-one element \(\gamma\) approximating
\(ghg^{-1}\) outside \(\ell\)
\cite[Main Theorem~28.5.3]{Voight}. In
\(gh=\gamma g(g^{-1}\gamma^{-1}gh)\), the final factor can be
taken in the right stabilizer of the vector away from \(\ell\).
At \(\ell\) it has norm one and acts trivially. Thus the
exceptional constituents are invariant under the entire finite
adelic norm-one group. On \(\mathcal K\) they depend only on the
joint norm squareclasses in \(G\). Indeed, at the odd level primes
a square norm can be scaled to one, and outside the level the
order units already act trivially. Conversely, each product of
local signs comes from a quadratic Dirichlet character composed
with reduced norm. This is projectively well-defined since
scalar norms are squares, and it has trivial real type since
real reduced norms are positive. By \eqref{eq:hecke} its
eigenvalue is \(\sigma(m)\) times its value on the norm class of
\(m\). These sign functions span the scalar exceptional space.

The remaining constituents have infinite-dimensional factors
at split places. Jacquet--Langlands gives cuspidal
\(\mathrm{GL}_2\) representations of holomorphic weight \(2k+2\)
and trivial central character, unramified at every \(\ell\mid m\)
\cite[Theorem~14.4 and Proposition~5.20]{JL}.
Let \(\alpha_\ell,\beta_\ell\) be their unitary Satake parameters.
Their product is one. The classical Hecke roots are
\(\ell^{k+1/2}\alpha_\ell,\ell^{k+1/2}\beta_\ell\), so Deligne's
bound gives \(|\alpha_\ell|=|\beta_\ell|=1\)
\cite[Theorems~4.9 and~5.1]{DeligneModular}
\cite[Theorem~1.6]{DeligneWeil}.

For the operator at \(\ell^a\) the all-ideal right cosets are
represented by upper triangular matrices with diagonal
\(\ell^i,\ell^{a-i}\) and upper right entry modulo \(\ell^i\),
for \(0\le i\le a\). In normalized unramified induction
\cite[Section~3]{JL}, the spherical vector has value
\(\ell^{a/2-i}\alpha_\ell^i\beta_\ell^{a-i}\) there.
Summing the \(\ell^i\) choices of the upper right entry gives the
eigenvalue
\[
 \ell^{a/2}\sum_{i=0}^{a}\alpha_\ell^i\beta_\ell^{a-i}.
\]
This includes imprimitive ideals and equals the classical
coefficient at \(\ell^a\) divided by \(\ell^{ak}\). In particular,
the recurrence is
\(T_{\ell^a}=T_\ell T_{\ell^{a-1}}-\ell T_{\ell^{a-2}}\) for
\(a\ge2\). The eigenvalue has modulus at most
\((a+1)\ell^{a/2}\). Multiplicativity now gives
\eqref{eq:hecke-norm} with the same bound at every weight and level.

\subsection{A finite test with a prescribed direction}
For a test \(b:G\to\C\) and a unit vector \(v\) in degree \(k\), define
\[
 F(g)=\frac1{|\Gamma|}\sum_{\gamma\in\Gamma}
                      \rho_k(\gamma)v\, b(\gamma^{-1}g).
\]
This is the orthogonal projection of \(b(g)v\) onto equivariant
functions in the full counting norm on \(G\). Its coset norm is at most
\(\|b\|_{\ell^2(G)}/\sqrt{|\Gamma|}\).
For another unit vector \(w\), reindexing the shell by the left
unit lift of each \(\gamma\) gives
\begin{equation}\label{eq:projection}
 12\langle w,(T_mF)(1)\rangle
 =\sum_{\alpha\in\mathcal Q_m}
       b(\alpha_{\rm fin})\,\langle w,\rho_k(u_\alpha)^{-1}v\rangle .
\end{equation}
Here the inner product is linear in its second variable.
The averaging factor cancels because for a lift \(\gamma_0\)
the substitution \(\alpha=\gamma_0\alpha'\) cancels
\(\rho_k(\gamma)\) against the archimedean multiplier.
This formula allows \(b\) to distinguish the two components.

Writing \(\alpha=z+yj\), put \(t_\alpha=(N(z)-N(y))/m\).
For \(v=w\) the unit vector proportional to \(X^kY^k\), the
matrix coefficient in \eqref{eq:projection} equals \(P_k(t_\alpha)\),
where \(P_k\) is the degree-\(k\) Legendre polynomial with \(P_k(1)=1\).
Indeed, the polynomial action gives the binomial formula
\[
 P_k(t_\alpha)=\sum_{j=0}^k(-1)^j\binom{k}{j}^2
               (N(z)/m)^{k-j}(N(y)/m)^j.
\]
Passing to \(u_\alpha^{-1}\) leaves the coefficient unchanged.

In degree zero, projecting onto the sign characters replaces \(b\)
by its average conditioned on the joint norm squareclasses of
\(m\). Indeed these characters form an orthogonal basis for
functions of the squareclasses, and the eigenvalue signs evaluate
them at \(m\). The averaging defining \(F\) preserves these means.
For a product of local tests the conditional mean is the product
of their conditional means. At each positive degree, evaluation
at one coset has norm at most one, so \eqref{eq:projection} is
bounded in absolute value by \(12\tau(m)\sqrt m\|b\|_{\ell^2(G)}\).
The same error bound applies to the scalar cuspidal part.

\subsection{The cost of a translated window}
Choose a fixed smooth \(0\le\phi\le1\) of compact support in
\((-1/2,1/2)\) and integral \(I_\phi>0\). Put \(h=H/n\),
\(b=M/n\), and \(\lambda=x/n\). We use the bump
\[
 W(t)=\phi\!\left(\frac{(1+t)/2-\lambda}{h}\right).
\]
As \(h\le b\), on its support \(1-t^2\asymp b\), and the
support has length at most \(2h\). The constant Legendre
coefficient of \(W\) is \(I_\phi h\).

\begin{lemma}\label{lem:bump}
The Legendre expansion \(W=\sum_{k\ge0}c_kP_k\) converges
uniformly and absolutely, with \(\sum_{k\ge1}|c_k|\ll_\phi\sqrt{b/h}\).
\end{lemma}
\begin{proof}
We use \(|P_k(t)|\ll k^{-1/2}(1-t^2)^{-1/4}\) for \(-1<t<1\).
For a short proof, let \(Y(u)=\sqrt{\sin u}\,P_k(\cos u)\) on
\(0<u<\pi\). The Legendre equation gives \(Y''+QY=0\) with
\(Q=(k+1/2)^2+1/(4\sin^2u)\).
The derivative of \(Y^2+(Y')^2/Q\) is
\(-Q'(Y')^2/Q^2\ge0\) up to \(\pi/2\). At \(\pi/2\) this energy
is \(O(k^{-1})\), by \(P_{2j}(0)=(-1)^j4^{-j}\binom{2j}{j}\),
\(P_{2j+1}(0)=0\) and \(P_k'(0)=kP_{k-1}(0)\).
Reflection gives the bound on both halves.

Write \(\mathcal L=-\frac d{dt}((1-t^2)\frac d{dt})\).
In the rescaled coordinate \(v=(t-2\lambda+1)/(2h)\), the operator
is \(b/h^2\) times a differential operator of order two with
coefficients and derivatives bounded on the support of \(\phi\).
Indeed its two coefficients in that scale are
\(-(1-t^2)/(4b)\) and \(th/b\). Integrating by parts twice
using \(\mathcal L P_k=k(k+1)P_k\), and also using the direct integral,
gives
\[
 |c_k|\ll_\phi h\sqrt{k}\,b^{-1/4}
       \min\!\left(1,\left(\frac{b}{h^2k^2}\right)^2\right).
\]
The factor \(\sqrt{k}\) includes the coefficient normalization
\((2k+1)/2\). Summing below and above \(K=\sqrt b/h\) gives
\(O_\phi(h b^{-1/4}K^{3/2})=O_\phi(\sqrt{b/h})\).
Since \(|P_k|\le1\) on \([-1,1]\), the series converges uniformly.
Its Legendre coefficients identify the sum with \(W\) by polynomial
density.
\end{proof}

Combining this lemma with \eqref{eq:projection}, a finite product
test of modulus at most one has weighted shell count
\begin{equation}\label{eq:angular}
 12\sigma(m)I_\phi h\prod_{p}\text{local conditional mean}
 +O_\phi\!\left(\sqrt{|G|}\,\tau(m)\sqrt m\sqrt{M/H}\right).
\end{equation}
The same \(W\) is used when the target is divided by a common factor.
Both component norms scale together and \(t_\alpha\) is unchanged.

\section{Local densities at all target valuations}\label{sec:local}
On a shell at \(m\), put \(A_1=N(z)\) and \(A_2=N(y)\).
The bump \(W\) restricts both norms to be positive. At a variable
prime \(p\ge5\) write \(g_m(p)\) for the density of the bad test
\(p\mid A_1A_2\). For \(p\nmid m\) this is the conditional mean in
\eqref{eq:angular}. At \(p\mid m\) we derive its analogue by reducing
to coprime Hecke indices.

\subsection{Coprime target and level}
At \(3\) the ordinary family requires no test, since \(m\equiv2\bmod3\)
forces both norms to be units modulo \(3\). In the mixed family
\(m\equiv1\bmod3\), and we require the prescribed component to have
\(3\)-adic valuation exactly one. Among the nine elements of \(R/3R\)
there are six units, two nonzero nonunits, and zero. The nonzero
nonunits have norm valuation one, since \(1-\omega\) has norm \(3\)
and its square is associated to \(3\). Norm one modulo \(3\)
requires one unit and one nonunit, giving \(2\cdot6\cdot3=36\)
raw quaternion residues. Exactly \(6\cdot2\) have a unit in one
prescribed component and a nonzero nonunit in the other. Passing to projective classes preserves the ratio,
so either mixed orientation has mean \(1/3\).

At an inert \(p\nmid m\), by anisotropy a bad point has one
component zero modulo \(pR\). Its image in
\(\PGL_2(\F_p)\) lies in the normalizer of the nonsplit torus
\((R/pR)^\times/\F_p^\times\). This is the union of that torus and
its translate by \(j\), with \(2(p+1)\) elements out of
\(p(p^2-1)\). The norms of the torus elements cover the squareclasses
equally, so \(g_m(p)=2/[p(p-1)]\) in either class.

At a split \(p\nmid m\), the two factors of the Eisenstein
subalgebra occupy the diagonal of \(M_2(\F_p)\). The other
Eisenstein component occupies the off-diagonal. Their norms are
the products of the respective entries, up to sign.
For a fixed nonzero determinant, the first row is uniform on
nonzero rows, each with \(p\) completions. Thus a given entry is
zero with probability \(1/(p+1)\). Two zero entries are possible
only when they form an opposite pair. For either such pair there
are \(p-1\) matrices of the given determinant, giving probability
\(1/[p(p+1)]\). Inclusion-exclusion gives
\begin{equation}\label{eq:B}
1-g_m(p)=B_p:=1-\frac4{p+1}+\frac2{p(p+1)}
             =\frac{(p-1)(p-2)}{p(p+1)}.
\end{equation}
For the exact condition \(v_p(A_i)=1\) at either one prescribed
component, exactly one of its two entries is zero modulo \(p\),
and this entry must be nonzero modulo \(p^2\). All matrix lifts of
an invertible residue are invertible, and determinant squareclasses
are determined modulo \(p\). The conditional probability of the
last lift condition is \((p-1)/p\). Therefore the mean of
\(v_p(A_i)=1\) at level \(p^2\) is
\begin{equation}\label{eq:t}
 t_p=\frac{2(p-1)^2}{p^2(p+1)},\qquad
 a_p:=\frac{t_p}{2B_p}=\frac{p-1}{p(p-2)} .
\end{equation}
These means hold on projective groups as well, since all tests
are unchanged by a unit scalar.

\subsection{Primes dividing the target}
For inert \(p\mid m\), divisibility of one norm by \(p\) forces
the other to be divisible by \(p\) as well, so both components
are in \(pR\). The bad set is empty if \(v_p(m)=1\).
Otherwise it is obtained by multiplying both coordinates by \(p\)
from the shell at \(m/p^2\). This preserves \(W\) and the other
prime tests. Hence for \(e=v_p(m)\) the scalar density is
\begin{equation}\label{eq:inert}
g_m(p)=0\ (e=1),\qquad
g_m(p)=\frac{\sigma(p^{e-2})}{\sigma(p^e)}\ (e\ge2).
\end{equation}

Now let \(p=\pi\bar\pi\) be split in \(R\), and fix \(r\) prime
to \(p\). At target \(p^e r\), let \(U_e\) be the weighted count
with no restriction at \(p\), and \(G_e\) the count with both
norms units at \(p\). Carry along \(W\) and the other prime
tests. At target \(r\), let \(L_0\) count with the additional
test
\[
 \frac12\bigl(\mathbf1_{\pi\mid z}+\mathbf1_{\bar\pi\mid z}
        +\mathbf1_{\pi\mid y}+\mathbf1_{\bar\pi\mid y}\bigr).
\]
It has mean \(2/(p+1)\) by the single-entry calculation.
If one norm has valuation \(v>0\), the \(v+1\) allocations of its
two prime factors are equally frequent. The average of the test
at that pair of norms is \(v/(v+1)\). The conditions at other
primes are unchanged on reallocating these factors. Indeed,
dividing out or inserting factors above \(p\) acts by units at
every other prime, preserving all the norm valuations and the
individual split-prime divisibilities used in the tests.

Classify a pair of norms on the shell \(p^e r\) by the minimum
\(j\) of its two \(p\)-adic valuations. If \(j<e\), both valuations
are \(j\); removing \(p^j\) from each norm gives a \(G_{e-j}\)
pair with allocation multiplier \((j+1)^2\). If \(j=e\), the
residual target is \(r\). When both residual norms are units
the multiplier is \((e+1)^2\); if one has valuation \(v>0\),
the ratio of allocation counts is
\((e+1)(v+e+1)/(v+1)=(e+1)^2-e(e+1)v/(v+1)\).
In either reduction the normalized norm coordinate stays fixed.
This proves
\begin{equation}\label{eq:recurrence}
 U_e=\sum_{j=0}^{e-1}(j+1)^2 G_{e-j}+(e+1)^2 U_0-e(e+1)L_0
 \qquad(e\ge1).
\end{equation}

The generating series for \((j+1)^2\) is
\(A(T)=(1+T)/(1-T)^3\), and that for \(e(e+1)\) is
\(C(T)=2T/(1-T)^3\). Thus
\[
 \sum_{e\ge1}G_eT^e=\frac{\sum_{e\ge0}U_eT^e}{A(T)}
                     -U_0+\frac{C(T)}{A(T)}L_0 .
\]
The coefficients of \(1/A(T)\) have absolute value at most eight,
and those of \(C(T)/A(T)\) at most two. The bad count \(U_e-G_e\)
is therefore a combination of unrestricted counts at lower
\(p\)-exponents and the \(L_0\) count at the unit target, with
coefficient absolute sum at most \(8e+2\). Applying the same
identity to the scalar prediction from \eqref{eq:angular} uses
\(U_e=\sigma(p^e)\), \(U_0=1\), and \(L_0=2/(p+1)\) after
factoring out the other local factors. Substitution gives
\(\sum_{e\ge1}G_e T^e=\frac{(p-1)^2}{p+1}\frac{T}{1-pT}\)
for those scalar values. In particular
\begin{equation}\label{eq:split}
 1-g_m(p)=\frac{(p-1)^2p^{e-1}}{(p+1)\sigma(p^e)}
                 \qquad(e=v_p(m)\ge1,\ p\text{ split}).
\end{equation}

\paragraph{Example at a split square.}\label{sec:split-example}
Take \(m=7^2r\) with \(7\nmid r\). The first two cases of
\eqref{eq:recurrence} give
\[
 U_2-G_2=4G_1+9U_0-6L_0=4U_1-7U_0+2L_0 .
\]
The term at \(7r\) is unrestricted at \(7\), and the only
remaining \(7\)-level test is \(L_0\) at \(r\). For the main
term, factor out the other primes and insert
\(U_2=57,\ U_1=8,\ U_0=1,\ L_0=1/4\).
The resulting main term for \(G_2\) is \(63/2\), giving
\(1-g_m(7)=(63/2)/57\).
Including the shell factor \(\sigma(7^2)/7^2=57/49\) gives
\(9/14\). This is \(B_7\cdot(6/5)\), since \(B_7=15/28\)
at a unit target. In the cofactor sum of Section~\ref{sec:cofactor-log}
this gives weight \(6/5\) without a common factor \(7\), and
\((6/5)/7\) with that common factor, since the residual target
\(7r\) still has positive valuation at \(7\). In general
\eqref{eq:split} gives \(E_p=(p-1)/(p-2)\) at every positive
valuation.

\subsection{The joint estimate}
Fix squarefree coprime \(u,v\) made of split primes coprime to
\(m\). Place the condition \(v_p(A_1)=1\) at \(p\mid u\) and
\(v_p(A_2)=1\) at \(p\mid v\). Give each coordinate pair weight
\(W(t)/\tau(uv)\). Here and below \(m\) is even and prime to \(3\).
Keep the odd-component condition and the prescribed condition
at \(3\). For any squarefree \(\ell\) coprime to \(6uv\), let
\(S_m(\ell)\) be the weighted count additionally bad at every
prime dividing \(\ell\), and put
\[
 X_m=I_\phi h R_o(m)\beta_m,\qquad
 \beta_m=\begin{cases}1,&m\equiv2\bmod6,\\1/3,&m\equiv4\bmod6.\end{cases}
\]
Extend \(g_m\) multiplicatively on squarefree products.

\begin{lemma}\label{lem:joint}
For \(m\le n\) as above and any fixed \(\zeta>0\),
\begin{equation}\label{eq:joint}
 S_m(\ell)=X_m\prod_{p\mid uv}(t_p/2)\,g_m(\ell)
   +O_{\phi,\zeta}\!\left((uv)^3\ell^{3/2}n^{1/2+\zeta}
                                           \sqrt{M/H}\right).
\end{equation}
\end{lemma}
\begin{proof}
First omit oddness. If the level is coprime to \(m\), apply
\eqref{eq:angular}. A factor at \(p\mid uv\) has size
\(p^3|\PGL_2(\F_p)|<p^6\) and a factor at \(p\mid\ell\)
has size less than \(p^3\). Thus the norm cost is
\(O((uv)^3\ell^{3/2})\), the fixed factor at \(3\) being bounded.

For \(p\mid(\ell,m)\), eliminate the test by inert scaling or
\eqref{eq:recurrence}. Each time, the remaining test at the
eliminated prime is present only at a reduced target coprime to it.
The allocation calculation works with any of the other tests
still imposed. In particular the test at a split prime produced
by an earlier elimination is an average of individual
divisibilities, all invariant under changes of allocation at
another prime. The coprime conditional means are independent of
the nonzero determinant class at those other primes. Extracting
split powers and inert squares does not change the target class
modulo \(3\), and the prescribed orientation is preserved. Hence
the scalar terms factor, with local factors
\eqref{eq:inert} and \eqref{eq:split}. The total coefficient variation
in the elimination is at most \(\tau(m)^4\), because
\(8e+2\le(e+1)^4\). Each resulting target divides \(m\), and the
remaining level divides the original level. The product of these
losses with the divisor factor in \eqref{eq:angular} is absorbed
by \(O_\zeta(n^\zeta)\).

Finally subtract the shell at \(m/4\) if it exists. Rescaling its
coordinates by \(2\) preserves the normalized angle, the odd-prime
tests and the orientation at \(3\). The valuations at all odd
primes agree with those for \(m\), so the scalar \(\sigma(m)\)
is replaced by \(\sigma(m)-\sigma(m/4)\). This gives
\eqref{eq:joint}.
\end{proof}

\section{The lower sieve and the split factors}\label{sec:sieve}
Fix \(\kappa,\xi>0\) with \(4\kappa+\frac52\xi<\varepsilon/2\).
We use a divisor level \(D=n^\xi\) and prime bounds \(Z=n^\eta\)
and \(Y=n^{\lambda_0}\), with \(0<\lambda_0<\eta\) chosen below.
Consider triples of pairwise coprime squarefree integers \(d,u,v\),
supported on split primes at most \(Y\), satisfying
\begin{equation}\label{eq:triples}
 d\mid n,\qquad (uv,n)=1,\qquad duv\le n^\kappa .
\end{equation}
On the shell \(m=n/d\) use the prescriptions of Lemma~\ref{lem:joint}
at \(u,v\), and require \(p\nmid A_1A_2\) at every other prime
\(5\le p\le Z\), including the primes of \(d\).
Since \(d\equiv1\bmod6\), the target \(m\) belongs to the same
family. The weight on the shell is
\(W(t)=\phi((dA_1-x)/H)\), so after rescaling the norm values by
\(d\) the original window is enforced.

\subsection{Finite inclusion-exclusion}
We give the sieve weights together with the level they require.
Partition the sieving primes into blocks
\(Z^{2^{-j-1}}<p\le Z^{2^{-j}}\) for \(j\ge0\).
For the main-term probabilities \(g_m(p)\), the sum in each block
is bounded by an absolute \(C\) and
\(v_j=\prod_{p\text{ in block }j}(1-g_m(p))\) is bounded below
by an absolute \(c>0\). In fact the formulas in
Section~\ref{sec:local} lie uniformly below one and are \(O(1/p)\).
Prime harmonic Mertens gives a bounded sum of \(1/p\) in each block.
Omitting primes, including the primes of \(uv\), preserves these
bounds. We also use the progression version of Mertens modulo \(3\),
\(\sum_{p\le T,\ p\equiv a\ (3)}1/p=\tfrac12\log\log T+O(1)\)
for \(a=1,2\); see \cite[Lemma~2.4]{LM}.

In block \(j\) truncate inclusion-exclusion at consecutive orders
\(k_j=2\lceil(K+j)/2\rceil\) and \(k_j-1\), where \(K\) is a fixed
large integer. Denote the resulting polynomials in the bad-prime
indicators by \(U_j,L_j\). Bonferroni gives
\(0\le\mathbf1_j\le U_j\) and \(L_j\le\mathbf1_j\), where
\(\mathbf1_j\) indicates avoidance in the block. Telescoping
\(\prod_j U_j-\prod_j\mathbf1_j\) gives the lower polynomial
\begin{equation}\label{eq:lower-poly}
 \prod_j\mathbf1_j\ge
 \prod_j U_j-\sum_j(U_j-L_j)\prod_{i\ne j}U_i .
\end{equation}
Under independent bad events with probabilities \(g_m(p)\),
\[
 \mathbf E U_j\ge v_j,\qquad
 \mathbf E(U_j-L_j)\le \frac{C^{k_j}}{k_j!}.
\]
Choose \(K\) so that the sum over \(j\) of \(C^{k_j}/(c k_j!)\)
is at most \(1/2\). The expectation of the lower polynomial is
at least \(\frac12\prod_j v_j\). Also
\(0\le \mathbf E U_j-v_j\le\mathbf E(U_j-L_j)\), so the expectation
of the upper polynomial \(\prod_j U_j\) is \(O(\prod_j v_j)\).
We will use this upper estimate later with integer-sieve
probabilities.

Put \(S=\sum_{j\ge0}k_j2^{-j}<\infty\) and \(\eta=\xi/S\).
Every divisor in either polynomial is at most \(Z^S=D\).
The upper coefficients have absolute value at most one, and
the lower coefficients at most one plus the number of nonempty
blocks, which is \(O(\log\log(3n))\). Apply \eqref{eq:joint}
to the lower polynomial with the nonnegative shell weights.
For each triple the main contribution is at least
\begin{equation}\label{eq:lower-main}
 \tfrac12 X_m V_m(Z)\prod_{p\mid uv}a_p,\qquad
 V_m(Z)=\prod_{5\le p\le Z}(1-g_m(p)).
\end{equation}
Here \(a_p=t_p/(2B_p)\) accounts for replacing the ordinary
sieve factor at a prime of \(uv\) by the exact-valuation factor.
Summing the remainders over all divisors uses
\(\sum_{\ell\le D}\ell^{3/2}\ll D^{5/2}\). There are at most
\(n^{\kappa+o(1)}\) triples because a fixed product \(duv\)
has at most \(\tau(duv)^2\) decompositions. The total absolute
remainder, for any \(\zeta>0\), is
\[
 O_{\phi,\zeta}\!\left(
       n^{1/2+4\kappa+(5/2)\xi+\zeta}\sqrt{M/H}\right)
   =o_\varepsilon(H/\log n).
\]
The equality uses
\(\sqrt{nM/H}\le Hn^{-\varepsilon/2}\)
and the strict choice of \(\kappa,\xi\), with \(\zeta\) small
enough. All of these estimates are uniform in \(x,H\).

\subsection{The logarithm gained from small factors}\label{sec:cofactor-log}
The target-independent product
\[
 V_*(Z)=
 \prod_{\substack{5\le p\le Z\\ p\ {\rm split}}}B_p
 \prod_{\substack{5\le p\le Z\\ p\ {\rm inert}}}b_p,
 \qquad b_p=1-\frac2{p(p-1)}
\]
satisfies \(V_*(Z)\asymp(\log Z)^{-2}\), since
\(B_p=1-4/p+O(p^{-2})\) and the inert factors have a convergent
positive product. To compare this with \(V_m(Z)\), include the
shell factor \(R_o(m)/m\) from \eqref{eq:odd-shell}.
For \(p^e\Vert m\) split and \(p\le Z\),
\[
 \frac{\sigma(p^e)}{p^e}(1-g_m(p))
       =\frac{(p-1)^2}{p(p+1)}=B_p E_p,\qquad
 E_p=\frac{p-1}{p-2}.
\]
For \(p^e\Vert m\) inert and \(p\le Z\) the combined factor
is \(1+1/p=b_p E_p\), by \eqref{eq:inert}, including \(e=1\).
The inert primes dividing \(m\) and \(n\) are the same.
For the inert primes dividing \(m\) beyond \(Z\), the shell
factor is at least \(1+1/p=b_pE_p\), and the product of the
extra factors \(b_p\) is bounded below. All split enhancements
are at least one. Using \(h m=(H/n)m=H/d\), we get
\begin{equation}\label{eq:target-gain}
 X_m V_m(Z)\gg_\phi
     \frac H d J(n)V_*(Z)
     \prod_{\substack{p\le Y,\ p\ {\rm split}\\ p\mid m}}E_p .
\end{equation}

Temporarily drop the product bound in \eqref{eq:triples}.
The sum of the cofactor weights in \eqref{eq:lower-main}
and \eqref{eq:target-gain} factors over split \(p\le Y\).
For \(p\nmid n\) the three options of absence, inclusion in
\(u\), or inclusion in \(v\) have weights \(1,a_p,a_p\).
If \(p^e\Vert n\), the two options of absence or inclusion in
\(d\) have weights \(E_p,E'_p/p\), where \(E'_p=1\) for \(e=1\)
and \(E'_p=E_p\) for \(e\ge2\).
The compensation at target divisors is the identity
\[
 E_p+\frac{E'_p}{p}\ge E_p+\frac1p=1+2a_p .
\]
Thus the full cofactor sum is at least
\(\prod_{p\le Y,\ p\ {\rm split}}(1+2a_p)\asymp\log Y\),
since \(a_p=1/p+O(p^{-2})\).

Normalize the local weights to independent probabilities.
At \(p\nmid n\) the probability of \(p\mid duv\) is
\(2a_p/(1+2a_p)\ll p^{-1}\), and at \(p\mid n\) it is at most
\(1/p\). Hence \(\mathbf E\log(duv)\le C'\log Y\), uniformly
in \(n\), by \(\sum_{p\le Y}(\log p)/p\ll\log Y\).
Take \(0<\lambda_0<\min(\eta,\kappa/(3C'))\). Markov's inequality
then retains at least half the weight on \(duv\le n^\kappa\).
Inserting \eqref{eq:target-gain} into \eqref{eq:lower-main}
and subtracting the aggregate remainder proves that the total
retained weight is
\begin{equation}\label{eq:weight}
 \gg_\varepsilon H J(n)\frac{\log Y}{(\log Z)^2}
       \gg_\varepsilon \frac{HJ(n)}{\log n}.
\end{equation}

\section{From coordinates to distinct pairs}\label{sec:count}
For each \(d\) fix \(\rho_d\in R\) of norm \(d\), obtained by
multiplying one factor above each split prime of \(d\).
Embed the residual shell at \(m=n/d\) in the original shell
by \((z,y)\mapsto(\rho_d z,\rho_d y)\). This preserves the
fractional norm coordinate, the prescribed orientation at \(3\),
and inert divisibility above \(Z\).

Each original value pair recovers its triple \(d,u,v\).
At split primes up to \(Y\), the only possible positive valuation
is one, occurring in both norms at the primes of \(d\), only
in the first at \(u\), and only in the second at \(v\).
The images for distinct triples are consequently disjoint.
This makes it possible to bound the remaining losses on the
original shell without summing an error over the triples.

\subsection{Large inert primes}
If an inert prime \(p>Z\) divides the first component norm,
that component is \(p\beta\). The norm of \(\beta\) belongs to
\([(x-H/2)/p^2,(x+H/2)/p^2]\). For the other component use
\(n-x\) in place of \(x\). Each resulting annulus has
squared-radius width \(H/p^2\) and outer radius at most
\(\sqrt n/p\). Covering each lattice point by a translated
fundamental cell, of fixed area and bounded radius, gives
\(O(H/p^2+\sqrt n/p+1)\) possible \(\beta\). Indeed the cells
lie in the annulus obtained by moving the two boundaries by
a fixed distance, truncating the inner radius at zero.
The complementary positive norm has at most
\(6\max_{a\le n}\tau(a)=n^{o(1)}\) coordinate choices.

Since all weights are at most one, summing on the original
shell over the two components and \(Z<p\le\sqrt n\) bounds
the lost weight by
\[
 n^{o(1)}\left(\frac H Z+\sqrt n(1+\log n)\right)
             =o_\varepsilon(H/\log n).
\]
Here \(Z\) is a fixed power of \(n\), and the hypothesis
\(H^3\ge n^{1+\varepsilon}M\) with \(M\ge H\) gives
\(H\ge n^{(1+\varepsilon)/2}\). Thus \eqref{eq:weight}
persists after removing every inert prime. Oddness and the
condition at \(3\) now give the required norm values.

\subsection{The weight of one value pair}
On the residual shell write \(A_1=uk_1\) and \(A_2=vk_2\).
All split prime factors of \(k_1,k_2\) exceed \(Z=n^\eta\).
Each of these two numbers has at most \(1/\eta\) split factors
counted with multiplicity, and the possible single factor \(3\)
has no effect on coordinate multiplicity.
The split primes in \(uv\) occur to valuation one. By the
multiplicity formula in Section~\ref{sec:angular}, the weighted
coordinate multiplicity, after division by \(\tau(uv)\),
is at most
\[
 36f(k_1)f(k_2)\le36\cdot 2^{2/\eta}.
\]
Each original pair determines its triple and residual pair.
Dividing the surviving weight by this constant proves the lower
bound of Theorem~\ref{thm:main}.

\subsection{The upper bound on integers}
For the upper bound, write \(s=s_1\) in the interval of length
\(H\). At every inert \(p\ge5\) we must exclude the residues
\(0,n\bmod p\). Let \(\nu(p)\) be their number, namely one if
\(p\mid n\) and two otherwise, and extend \(\nu\)
multiplicatively to squarefree products of these primes.
The Chinese remainder theorem and the count in a residue
class give \(H\nu(\ell)/\ell+O(\nu(\ell))\) integers in
all the excluded sets for \(p\mid\ell\), uniformly in the
interval location and endpoint convention.

Use the upper polynomials of Section~\ref{sec:sieve} with
probabilities \(\nu(p)/p\), sieving to \(Z'=n^{c_1}\) for a
sufficiently small fixed \(c_1>0\).
The block probability bounds hold as before. Choose \(c_1\)
so that the corresponding divisor level is \(D'\le n^{1/4}\).
The total error is \(O(D'\log(2D'))=o(H/\log n)\), since
\(\nu(\ell)\le\tau(\ell)\)
and \(\sum_{\ell\le D'}\tau(\ell)\ll D'\log(2D')\).
The upper main term is at most a constant times \(H\) times
\[
 \prod_{\substack{5\le p\le Z'\\ p\ {\rm inert}}}(1-2/p)
 \prod_{\substack{5\le p\le Z',\ p\mid n\\ p\ {\rm inert}}}
                        \frac{p-1}{p-2}
 \ll \frac{J(n)}{\log n}.
\]
This proves the upper bound and completes Theorem~\ref{thm:main}.
The argument for this upper bound in fact works in any interval
whose length is at least a fixed positive multiple of \(\sqrt n\).

\begin{proof}[Proof of Corollary~\ref{cor:ranges}]
In a fixed interval of shares, take a smaller fixed interval
of length at most the distance of its center from the nearest
endpoint of \([0,1]\). This satisfies the hypotheses with
\(H\asymp n\) for large \(n\), giving the lower bound. Use the
upper sieve on the full interval and on \([0,n]\) for the
unrestricted count.
At \(x=\lambda n\) with \(0<\lambda<1\), a width \(H=n^\theta\)
satisfies the condition for any fixed
\(0<\varepsilon<3\theta-2\).
Finally, inside \([a n^\theta,b n^\theta]\) take a subinterval
of width comparable to \(n^\theta\) and at most its midpoint.
Its distance to the nearer endpoint of \([0,n]\) is comparable
to \(n^\theta\), so one can use \(0<\varepsilon<2\theta-1\).
The upper sieve on the full interval gives the matching bound,
also for \(a=0\).
\end{proof}

\paragraph{AI assistance and review.}
ChatGPT (OpenAI) assisted with mathematical arguments, literature checks,
exposition, proof review and verification code.
\begin{thebibliography}{99}\small
\bibitem{Euler} W. Howe, G. Huffman, J. Kohlmeyer, and L. Kruer,
\emph{Euler's 1747 conjecture on sums of integers with restricted prime
factors}, Hexagon 2611.00370v1, 12 November 2026.
\url{https://hexagonmath.org/2611.00370v1}.
\bibitem{DietmannElsholtz} R. Dietmann and C. Elsholtz,
\emph{Sums of two numbers having only prime factors congruent to one
modulo four}, Banach Center Publ. 118 (2019), 177--179.
\url{https://doi.org/10.4064/bc118-11}.
\bibitem{Voight} J. Voight, \emph{Quaternion Algebras}, Graduate Texts in
Mathematics 288, Springer, 2021, Chapters 28 and 41.
\url{https://doi.org/10.1007/978-3-030-56694-4}.
\bibitem{RV} O. Ramar\'e and G. K. Viswanadham,
\emph{The trigonometric polynomial on sums of two squares, an additive
problem and generalisation}, arXiv:2509.23260v1, 2025.
\url{https://arxiv.org/abs/2509.23260v1}.
\bibitem{JL} H. Jacquet and R. P. Langlands,
\emph{Automorphic Forms on $\mathrm{GL}(2)$}, Lecture Notes in Mathematics
114, Springer, 1970, Theorem~14.4 and Section~5.
\url{https://publications.ias.edu/sites/default/files/automorphic-forms-on-gl2_rpl_9.pdf}.
\bibitem{DeligneModular} P. Deligne,
\emph{Formes modulaires et repr\'esentations $\ell$-adiques},
S\'eminaire Bourbaki, expos\'e 355 (1968/69), published 1971, 139--172.
\url{https://www.numdam.org/item/SB_1968-1969__11__139_0.pdf}.
\bibitem{DeligneWeil} P. Deligne, \emph{La conjecture de Weil. I},
Publications Math\'ematiques de l'IH\'ES 43 (1974), 273--307.
\url{https://www.numdam.org/item/PMIHES_1974__43__273_0.pdf}.
\bibitem{GV} M. Greenberg and J. Voight,
\emph{Lattice methods for algebraic modular forms on classical groups},
in Computations with Modular Forms, Contrib. Math. Comput. Sci. 6,
Springer, 2014, 147--179.
\url{https://jvoight.github.io/articles/orthohecke-240123.pdf}.
\bibitem{LM} D. Loughran and L. Matthiesen,
\emph{Frobenian multiplicative functions and rational points in fibrations},
J. Eur. Math. Soc. 26 (2024), 4779--4830.
\url{https://doi.org/10.4171/jems/1374}.
\bibitem{IwaniecHalf} H. Iwaniec, \emph{The half dimensional sieve},
Acta Arith. 29 (1976), 69--95.
\url{https://doi.org/10.4064/aa-29-1-69-95}.
\bibitem{Greaves} G. Greaves, \emph{On the representation of a number
in the form \(x^2+y^2+p^2+q^2\) where \(p,q\) are odd primes},
Acta Arith. 29 (1976), 257--274.
\url{https://doi.org/10.4064/aa-29-3-257-274}.
\end{thebibliography}

\end{document}
